\section{Conclusion}

We investigated how mixed and adaptive resource allocation strategies influence the exploration--exploitation tradeoff in batch multi-objective Bayesian optimization under evolving resource constraints in autonomous materials discovery campaigns. Fixed, rule-based adaptive, and LLM-guided adaptive policies were evaluated across a three-objective soft-magnetic alloy design problem ($M_s/H_c/H_V$ in the Fe--Co--Ni--V--Mo--C--Nb--Ti--Si system) and a two-objective melting-point/density problem ($T_m/\rho$ in the Ti--V--Nb--Mo--Hf--Ta--W system), spanning eight experimental configurations covering baseline and single- and dual-event conditions.
The central quantitative result is that the percentage gain in cumulative mutual information achieved by mixed allocation strategies substantially exceeded the corresponding percentage loss in hypervolume and AUC-HV relative to qEHVI. 

In the RHEA baseline campaign, Adaptive-$\beta$ accumulated approximately 37\% more cumulative MI than qEHVI ($19.54 \pm 0.41$ versus $14.23 \pm 0.45$ nats) while finishing within 0.2\% of qEHVI in final hypervolume and within 0.8\% in AUC-HV. In the Fe-Co-Ni magnet problem, where all strategies produced statistically indistinguishable final hypervolumes and AUC-HV values, Adaptive-$\beta$ accumulated approximately 49\% more MI than qEHVI at no measurable cost in optimization performance. Rule-Swap and Fixed Mixed fell between these values in the same ordering across both systems. 

The implication is that exploration-biased campaigns acquired a substantially more informative characterization of the Pareto-front neighborhood even when optimization convergence appeared uniform, a distinction whose practical value lies in the downstream utility of the experimental record for objective refinement and follow-up work.

Among the mixed allocation strategies, the adaptive policies outperformed Fixed Mixed in responsiveness to both evolving campaign statistics and mid-campaign resource events, advancing their normal allocation trajectory in response to constraints rather than altering the underlying strategy itself, while Fixed Mixed's predetermined exploration schedule could not be revised to account for the shortened campaigns.

The primary limitation of the current framework is that both LLM-based strategies operated exclusively on optimization statistics, drawing on hypervolume history, cumulative MI, and remaining campaign resources without access to any prior knowledge about the materials system being studied. This restriction enforced evidence-based decisions and avoided problem-specific bias, but it also meant that allocation decisions in the early campaign, before sufficient optimization history had accumulated, were made without knowledge of how composition relates to the objectives being optimized.

These findings point toward several directions for improving LLM-guided allocation agents in self-driving laboratory settings. Resource allocation agents equipped with prior information about a materials system could influence the selection of candidate points beyond simply choosing among acquisition function suggestions, directing early-campaign sampling toward regions of composition space that prior knowledge identifies as more likely to yield Pareto-relevant compositions. 

A further open question is the extent to which the mutual information advantages documented across all experiments translate into measurable improvements in Pareto-front prediction accuracy. Quantifying this relationship would provide a principled basis for selecting the exploration intensity that maximizes the long-term utility of an experimental campaign rather than optimizing solely for hypervolume within a fixed horizon. More broadly, this work establishes a generalizable framework for integrating probabilistic optimization with adaptive, language-driven decision-making, and demonstrates that LLM agents can play a structured and beneficial role as high-level controllers in multi-objective experimental campaigns operating under realistic, evolving constraints.